%%
%% UTF-8       UTF-8      UTF-8       UTF-8       UTF-8       UTF-8
%% UTF-8       UTF-8      UTF-8       UTF-8       UTF-8       UTF-8
%%
%% XeLaTeX  XeLaTeX  XeLaTeX  XeLaTeX  XeLaTeX  XeLaTeX  XeLaTeX  XeLaTeX  

% !TeX program = xelatex
\lccode`\-=`\-
%\defaulthyphenchar=127
\documentclass[12pt,a4paper,oneside]{extarticle}
%\usepackage{xunicode} % https://ctan.org/pkg/xunicode
%\usepackage{xltxtra} % https://ctan.org/pkg/xltxtra
\usepackage{fontspec} % https://ctan.org/pkg/fontspec
\usepackage{polyglossia} % https://ctan.org/pkg/polyglossia
\setdefaultlanguage{portuges}
%\setdefaultlanguage{portuguese}
%\setdefaultlanguage[spelling=old]{portuguese}
\usepackage{csquotes}

%\usepackage{ebgaramond}%[RawFeature={+onum,+pnum,+ss06}]
\usepackage[cmintegrals,cmbraces]{newtxmath}
%\usepackage{ebgaramond-maths}

\usepackage[no-math]{fontspec}
%\setmainfont{EBGaramond}[
%  Extension=.otf,
%  UprightFont=*-Regular,
%  ItalicFont=*-Italic,
%  BoldFont=*-Bold,
%  BoldItalicFont=*-BoldItalic,
%  Numbers={Proportional,OldStyle},
%  Ligatures=TeX
%]

%\setmainfont{TeX Gyre Pagella}[
%  Extension = .otf,
%  UprightFont = *-Regular,
%  ItalicFont = *-Italic,
%  BoldFont = *-Bold,
%  BoldItalicFont = *-BoldItalic,
%  Ligatures = TeX
%]



\clubpenalty=100
\linepenalty=100
\widowpenalty=100

%estética
% \usepackage[hyphens]{url}
\usepackage{xurl}
\urlstyle{same} % URLs na mesma fonte do texto (Times New Roman)
\AtBeginDocument{%
  \setlength{\Urlmuskip}{0mu plus 1mu} % Ajusta espaçamento entre caracteres em URLs
  \setlength{\emergencystretch}{3em} % Permite maior flexibilidade na justificação
  \def\UrlBreakPenalty{100} % Penalidade baixa para facilitar quebras
  \def\UrlBigBreakPenalty{50} % Penalidade para quebras em caracteres como /
  \sloppy % Relaxa a justificação para evitar ultrapassagem
}

\usepackage[hang,small,bf]{caption}
\usepackage[a4paper,left=2.5cm,right=2.5cm,top=2.5cm,bottom=2.5cm,includehead,includefoot]{geometry}
\linespread{1.3}

\usepackage[
  xetex,                      % Driver para XeLaTeX
  unicode=true,               % Suporte a caracteres Unicode (ex.: ç, ã)
  bookmarks=true,             % Cria bookmarks no PDF
  bookmarksopen=true,         % Bookmarks abertos por padrão
  bookmarksnumbered=true,     % Numera os bookmarks (ex.: 1, 1.1, 1.2)
  linktoc=all,                % Links clicáveis em sumário, figuras e citações
  breaklinks=true,            % Permite quebra de links longos
  colorlinks=true,            % Links com cor em vez de caixas
  citecolor=black,            % Cor das citações (preto)
  filecolor=black,            % Cor dos links de arquivos (preto)
  linkcolor=black,            % Cor dos links internos (preto)
  urlcolor=blue,              % URLs em azul
  pdfpagelabels=true,         % Adiciona rótulos às páginas do PDF
  plainpages=false,           % Diferencia numeração de páginas
  pdftitle={Desafios 12 ano},
  pdfauthor={Desafios 12 ano}, % Nome do autor
  pdfsubject={Desafios 12 ano},
  pdfkeywords={Desafios 12 ano},
  pdfcreator={XeLaTeX com hyperref},
  pdfproducer={XeLaTeX}
]{hyperref}

%% Definir um novo estilo de url
\makeatletter
\def\url@luisstyle{%
  \@ifundefined{selectfont}{\def\UrlFont{\sf}}{\def\UrlFont{\small\ttfamily}}}
\makeatother
\urlstyle{luis}

\setlength{\captionmargin}{1.5in}
\usepackage{fancyhdr}

\renewcommand{\thefootnote}{\alph{footnote}} %modificar as letras de notas de rodapé%
\makeatletter
\renewcommand\@makefnmark{\@textsuperscript{\normalfont(\@thefnmark)}}
\makeatother
\renewcommand{\fboxrule}{2pt} %para alterar a expessura das caixas
%\renewcommand{\fboxsep}{3ex} %para alterar a separação nas caixas

%para alterar a fonte dos items http://cnlart.web.cern.ch/cnlart/214/node68.html
%\renewcommand{\labelenumi}{\textbf{\arabic{enumi}.}}
\renewcommand{\labelenumi}{\textbf{\alph{enumi}.}}
%\renewcommand{\labelenumii}{\textbf{\alph{enumii}.}}
%\renewcommand{\labelenumii}{\textbf{(\alph{enumii})}}
\renewcommand{\labelenumii}{\textbf{(\arabic{enumii})}}

\long\def\symbolfootnote[#1]#2{\begingroup\def\thefootnote{\fnsymbol{footnote}}\footnote[#1]{#2}\endgroup}

%%%%%%%%%%%%%%%%%%%%%%%%%%%%
%%%%%%%%%%%%%%%%%%%%%%%%%%%%
%%%%%%%%%%%%%%%%%%%%%%%%%%%%
\usepackage{setspace}
%\singlespacing
%\onehalfspacing
%\doublespacing

\usepackage{comment}
\usepackage{siunitx}
%%%%%%%%%%%%%%%%%%%%%%%%%%%%
\usepackage[x11names]{xcolor} 
%%%%%%%%%%%%%%%%%%%%%%%%%%%%
\usepackage{relsize}
%\newcommand\CC{C\nolinebreak\hspace{-.05em}\raisebox{.4ex}{\relsize{-3}{\textbf{+}}}\nolinebreak\hspace{-.10em}\raisebox{.4ex}{\relsize{-3}{\textbf{+}}}}
\newcommand\CC{C\nolinebreak[4]\hspace{-.05em}\raisebox{.4ex}{\relsize{-3}{\textbf{++}}}}

%\usepackage{showframe}
\usepackage{pgf,tikz}       % - Elementos gráficos
\usepackage{tikzpagenodes}
\usetikzlibrary{backgrounds,calc,positioning}
\usetikzlibrary{positioning, shapes}

\pgfdeclarelayer{background}
\pgfdeclarelayer{foreground}
\pgfsetlayers{background,main,foreground}

\usepackage[
    backend=biber,
    style=numeric-comp,
    language=auto,
    natbib=false,
    block=space,
    isbn=false,
    url=true, 
    doi=true,
    mcite=true,
    eprint=false,
    sorting=none,
    maxnames=6
]{biblatex} % https://ctan.org/pkg/biblatex
\addbibresource{bibliografia.bib}
 % não é para alterar

\begin{document}
%espaço negativo para o cabeçalho
\vspace*{-2cm}
\begin{tikzpicture}[remember picture, overlay]
  \draw[line width = 0.5pt, color=black] ($(current page.north west) + (2.5cm,-2cm)$) -- ($(current page.north east) + (-2.5cm,-2cm)$);
\end{tikzpicture}
\vspace*{2mm}
\begin{minipage}[c][2.5cm][t]{.25\textwidth}
    \includegraphics[width=\linewidth]{logos/aeB06.pdf}
\end{minipage}\hfill
\begin{minipage}[c][2.5cm][t]{.7\textwidth}
	\begin{center}\Large\textbf{Agrupamento de Escolas de Benfica}\end{center}
	\vspace{-8.5mm}
	\begin{center}\large\textbf{Escola Secundária José Gomes Ferreira}\end{center}
	\vspace{-8.5mm}
	\begin{center}\large{Grupo Disciplinar 510}\end{center}
\end{minipage}
\begin{tikzpicture}[remember picture, overlay]
\draw[line width = 0.5pt, color=black] ($(current page.north west) + (2.5cm,-4.8cm)$) -- ($(current page.north east) + (-2.5cm,-4.8cm)$);
\end{tikzpicture}
\vspace{10mm}
\begin{center}
\Large{\textbf{Análise e Implementação Computacional da Geometria dos Favos de Mel}}\par 
\large{\textbf{Adrian Dias}}\par 
\large{\textbf{Nº 1, T 12º 3ª }}\par 
\Large{-- Disciplina de Física --}\par
\textbf{2025 - 2026}\\
\end{center}

\begin{tikzpicture}[remember picture, overlay]
  \draw[line width = 0.5pt, color=black] ($(current page.south west) + (2.5cm,3cm)$) -- ($(current page.south east) + (-2.5cm,3cm)$);
\end{tikzpicture}
 % pode alterar o título, o nome o número etc.

\begin{abstract}
Este relatório apresenta o desenvolvimento de um modelo matemático e a respectiva implementação computacional em \texttt{Python}/\texttt{turtle}, com vista à representação de uma forma geométrica presente na natureza. Seleccionou-se como objeto de estudo a estrutura hexagonal dos favos de mel, dada a sua relevância matemática e propriedades de optimização. São descritos os fundamentos teóricos, a metodologia de implementação e os principais resultados obtidos, com particular ênfase na modelação matemática subjacente.
\end{abstract}

\section{Introdução}
\subsection{Enquadramento e Objectivos}
O presente trabalho tem como objectivo principal a modelação matemática e implementação computacional de formas naturais, com particular enfoque nos padrões hexagonais observáveis em favos de mel. A selecção desta forma específica justifica-se pela sua fundamentação matemática robusta e pelas propriedades de optimização que exibe na natureza.

A abordagem adoptada integra conceitos de geometria euclidiana, teoria de tesselações e algoritmos computacionais, visando não apenas a reprodução gráfica da forma, mas também a compreensão dos princípios matemáticos que governam a sua estrutura.

\subsection{Fundamentos Matemáticos dos Favos de Mel}

\subsubsection{Conjectura do Favo de Mel e Princípios de Optimização}
A estrutura hexagonal dos favos de mel materializa a \textbf{Conjectura do Favo de Mel} (Honeycomb Conjecture), demonstrada matematicamente por Thomas Hales em 1999. Este teorema estabelece que, entre todas as partições do plano em regiões de área igual, a configuração hexagonal regular minimiza o perímetro total. A razão óptima perímetro-área é dada por:

\begin{equation}
\min \frac{P}{A} = \sqrt[4]{12} \approx 1,8612
\label{eq:honeycomb_conjecture}
\end{equation}

onde $P$ representa o perímetro e $A$ a área. Para o hexágono regular com lado $s$, obtém-se:

\begin{equation}
\frac{P}{A} = \frac{6s}{\frac{3\sqrt{3}}{2}s^2} = \frac{4}{\sqrt{3}s}
\label{eq:hexagon_ratio}
\end{equation}

\subsubsection{Tesselações Regulares do Plano}
A geometria euclidiana impõe restrições rigorosas às tesselações regulares do plano. Conforme demonstrado matematicamente, apenas três polígonos regulares preenchem completamente o espaço bidimensional sem sobreposições ou espaços vazios:

\begin{itemize}
    \item \textbf{Triângulos equiláteros} (notação de Schläfli: $\{3,6\}$)
    \item \textbf{Quadrados} (notação de Schläfli: $\{4,4\}$) 
    \item \textbf{Hexágonos regulares} (notação de Schläfli: $\{6,3\}$)
\end{itemize}

A condição matemática necessária e suficiente para tesselação regular é expressa por:

\begin{equation}
\frac{1}{m} + \frac{1}{n} = \frac{1}{2}
\label{eq:tessellation_condition}
\end{equation}

onde $m$ e $n$ são inteiros que satisfazem a equação para os três casos mencionados.

\subsubsection{Limitações ao Refinamento Hierárquico}
Os hexágonos regulares apresentam constrangimentos matemáticos inerentes no que concerne a sistemas de múltiplas escalas. A impossibilidade de subdivisão hierárquica perfeita decorre de três axiomas fundamentais:

\begin{enumerate}
    \item \textbf{Conservação de área}: $A_{\text{parental}} = \sum A_{\text{filhas}}$
    \item \textbf{Hierarquia simples}: Apenas um elemento parental por célula
    \item \textbf{Cobertura perfeita}: Ausência de espaços vazios ou sobreposições
\end{enumerate}

Estes princípios matemáticos explicam a preferência evolutiva por favos de dimensão uniforme em detrimento de sistemas com múltiplas escalas.

\section{Modelo Computacional Implementado}

\subsection{Arquitetura do Sistema de Coordenadas}
No script \texttt{Python/turtle}, a estrutura hexagonal é implementada através de um sistema de coordenadas baseado em vectores de rede hexagonal. A posição de cada hexágono é determinada por:

\begin{equation}
\vec{r}_{ij} = i \cdot \vec{a}_1 + j \cdot \vec{a}_2
\label{eq:hexagonal_grid}
\end{equation}

onde os vectores base da rede hexagonal são definidos como:

\begin{align}
\vec{a}_1 &= \left(1.5s, 0\right) \\
\vec{a}_2 &= \left(0.75s, \frac{\sqrt{3}}{2}s\right)
\label{eq:base_vectors}
\end{align}

sendo $s$ o comprimento do lado do hexágono.

\subsection{Especificações e Parâmetros do Modelo}
Para a implementação computacional, foram estabelecidos os seguintes parâmetros matemáticos:

\begin{itemize}
    \item \textbf{Dimensão do hexágono}: $s = 25$ unidades turtle
    \item \textbf{Número de camadas}: $n = 7$ camadas concêntricas
    \item \textbf{Total de hexágonos}: $N = 1 + \sum_{k=1}^{7} 6k = 127$ células
    \item \textbf{Ângulo de rotação}: $\theta = 60^\circ$ entre lados consecutivos
    \item \textbf{Área individual}: $A_{\text{hex}} = \frac{3\sqrt{3}}{2}s^2$
\end{itemize}

\subsection{Algoritmo de Construção Progressiva}
O modelo implementado simula o processo de construção natural através de um algoritmo de expansão radial, iniciando em pontos centrais e expandindo concentricamente. A progressão obedece à sequência:

\begin{equation}
N_k = 6k \quad \text{para} \quad k = 1, 2, \ldots, n
\label{eq:progressive_growth}
\end{equation}

onde $N_k$ representa o número de hexágonos na camada $k$-ésima, garantindo crescimento simétrico e geometricamente consistente.

\subsection{Implementação Computacional}

\begin{lstlisting}[language=python]
import turtle
import math

def draw_hexagon(t, size, depth=0.3):
    """Desenha um hexágono com Turtle, simulando profundidade."""
    t.pensize(2)  # Bordas mais grossas como na imagem
    for _ in range(6):
        t.forward(size)
        t.left(60)
    
    # Simula profundidade desenhando um hexágono interno menor
    t.penup()
    t.forward(size * depth)  # Move para dentro
    t.left(30)  # Ajusta para alinhar com o centro
    t.pendown()
    t.pensize(1)  # Bordas internas mais finas
    for _ in range(6):
        t.forward(size * (1 - depth))
        t.left(60)
    t.penup()
    t.goto(t.xcor() - size * depth, t.ycor())  # Volta ao ponto inicial
    t.setheading(t.heading() - 30)  # Corrige orientação

def generate_honeycomb(rows=8, cols=10, hex_size=20):
    """
    Gera uma estrutura de favos de mel hexagonal usando Turtle.
    
    Parâmetros:
    - rows: Número de linhas de hexágonos.
    - cols: Número de colunas de hexágonos.
    - hex_size: Tamanho do lado do hexágono.
    
    Baseado na tesselação {6,3} (Schläfli) e otimizado para parecer com a imagem.
    """
    t = turtle.Turtle()
    t.speed(0)  # Máxima velocidade
    t.hideturtle()  # Esconde o Turtle para melhor visualização
    
    # Configurações de cor baseadas na imagem (tons de mel)
    t.fillcolor("#FFD700")  # Dourado mel
    t.pencolor("#DAA520")  # Bege acastanhado para bordas
    
    # Distâncias para posicionamento hexagonal
    hex_width = math.sqrt(3) * hex_size
    hex_height = 2 * hex_size
    
    for row in range(rows):
        y = row * (3/2 * hex_size)  # Offset vertical
        for col in range(cols):
            # Offset horizontal para linhas pares/ímpar
            x_offset = (col * hex_width) + (row % 2) * (hex_width / 2)
            
            # Move para a posição inicial
            t.penup()
            t.goto(x_offset - (cols * hex_width / 2), -y - (rows * hex_height / 4))  # Centraliza
            t.pendown()
            
            # Desenha e preenche o hexágono
            t.begin_fill()
            draw_hexagon(t, hex_size, depth=0.3)
            t.end_fill()
    
    # Mantém a janela aberta
    turtle.done()

# Exemplo de uso
if __name__ == "__main__":
    generate_honeycomb(rows=8, cols=10, hex_size=20)
\end{lstlisting}


\section{Parte Experimental}

\subsection{Contexto e Abordagem Metodológica}

A componente experimental deste trabalho corresponde à elaboração e explicação do código para análise computacional da geometria hexagonal em favos de mel. Tal como num procedimento laboratorial, importa detalhar a lógica implementada, os algoritmos utilizados e as opções tomadas em cada etapa, de forma a permitir a replicação do processo.

A investigação parte da \textbf{Honeycomb Conjecture}, que estabelece o padrão hexagonal como a forma mais eficiente para particionar o plano em regiões de igual área com menor perímetro total. Esta fundamentação teórica justifica a escolha do hexágono regular como objeto central de estudo.

\subsection{Implementação Computacional}

A implementação foi realizada utilizando a biblioteca \texttt{turtle} do Python, escolhida pela sua capacidade de produzir representações gráficas bidimensionais de maneira simples e eficaz. O desenvolvimento seguiu uma abordagem baseada em tesselações \cite{mathstackexchange_tiling_notations_2022}, considerando notações padronizadas para empacotamento de polígonos.

\subsubsection{Algoritmo Principal}

O algoritmo central, implementado na função \texttt{generate\_honeycomb}, organiza uma estrutura hexagonal que reflete os padrões naturais dos favos de mel. A lógica de posicionamento utiliza um sistema de coordenadas baseado em deslocamentos regulares, assegurando uma disposição simétrica e contínua das células:

\begin{equation}
x_{\text{offset}} = \sqrt{3} \times \text{hex\_size} \times \text{col}
\end{equation}

\begin{equation}
y_{\text{offset}} = 1.5 \times \text{hex\_size} \times \text{row}
\end{equation}

Esta abordagem considera as propriedades de subpavimentação com hexágonos \cite{mathstackexchange_hexagon_subpaving_2023}, garantindo cobertura total do plano sem sobreposições.

\subsubsection{Desenho de Hexágonos Individuais}

A função \texttt{draw\_hexagon} foi desenvolvida para traçar cada hexágono individual, incorporando um elemento interno menor com um fator de profundidade de $0,3$, visando simular a tridimensionalidade observada na natureza. A generalização para outros polígonos regulares foi considerada com base em estudos sobre tesselações com pentágonos \cite{mathstackexchange_pentagon_tiling_2023}.

\subsection{Parâmetros e Configurações}

Foram estabelecidos parâmetros específicos para otimização da visualização:

\begin{itemize}
    \item Número de linhas: $\mathrm{rows} = 8$
    \item Número de colunas: $\mathrm{cols} = 10$  
    \item Tamanho do hexágono: $\mathrm{hex\_size} = 20$
    \item Cores: Preenchimento dourado (\textcolor[HTML]{FFD700}{\texttt{\#FFD700}}) e bordas em bege (\textcolor[HTML]{DAA520}{\texttt{\#DAA520}})
    \item Espessura das bordas: $\mathrm{pensize} = 2$
\end{itemize}

\subsection{Processo de Construção e Validação}

O processo de construção segue uma abordagem iterativa, simulando um crescimento progressivo a partir de um ponto central. A centralização da grade é assegurada por deslocamentos proporcionais às dimensões totais da estrutura:

\begin{equation}
x_{\text{start}} = -\frac{\text{cols} \times \sqrt{3} \times \text{hex\_size}}{2}
\end{equation}

\begin{equation}
y_{\text{start}} = -\frac{\text{rows} \times 1.5 \times \text{hex\_size}}{2}
\end{equation}

A validação geométrica considerou também aspectos de tesselações hiperbólicas \cite{mathstackexchange_hyperbolic_tessellations_2013}, embora o foco tenha permanecido no plano euclidiano.

\subsection{Replicabilidade e Extensões}

Esta metodologia permite que o código seja replicado em qualquer ambiente Python com a biblioteca \texttt{turtle} devidamente instalada. A arquitetura modular do sistema permite futuras extensões para:
\begin{itemize}
    \item Análise de eficiência de empacotamento
    \item Simulação de deformações estruturais
    \item Estudo de transições para outros polígonos regulares
    \item Aplicação em contextos de geometria hiperbólica
\end{itemize}

A implementação serve assim como ferramenta versátil para exploração computacional de padrões geométricos na natureza.









\section{Discussão dos Resultados}
Apresentam-se e analisam-se, nesta secção, as imagens geradas automaticamente pelo código.  
Não foram utilizadas capturas de ecrã, mas sim exportações directas produzidas pelo programa.  
Discutem-se as semelhanças e diferenças entre os resultados e a imagem de referência, identificando as causas dos desvios e avaliando a qualidade da aproximação obtida.

\section{Conclusões}
As conclusões são redigidas a partir da análise dos resultados.  
Evitam-se afirmações superficiais ou subjectivas; privilegiam-se observações fundamentadas, como, por exemplo:  
\begin{itemize}
    \item O modelo reproduz com fidelidade parcial a forma natural seleccionada.
    \item As limitações decorrem de aproximações matemáticas ou restrições do ambiente de programação.
    \item Futuras melhorias poderão incluir optimizações algorítmicas ou refinamentos gráficos.
\end{itemize}

\section{Bibliografia}
\printbibliography

\label{fim}
\end{document}
